\section{Related Work}
\label{sec:background}

\subsection{Spider Benchmarks}

Spider~1.0~\citep{yu2018spider} comprises 10{,}181 question-SQL pairs across
200 databases of varying complexity (simple to extra-hard).
The development set contains 1{,}034 questions spanning 20 databases and
serves as the standard evaluation split.
Execution accuracy (EX)---the fraction of questions for which the predicted
SQL returns the same result set as the gold SQL---is the primary metric.

Spider~2.0-Lite~\citep{lei2025spider2} is a substantially harder benchmark
targeting enterprise-scale analytics.
It comprises 547 questions; the local SQLite subset contains 135 questions
across 30 databases, featuring complex constructs such as window functions,
recursive CTEs, and multi-table analytical joins.
Most published systems achieve EX\,=\,0.00 on this benchmark without task-specific
adaptation.

\subsection{Text-to-SQL Approaches}

Two broad paradigms dominate the literature.
Fine-tuned approaches such as RESDSQL~\citep{li2023resdsql} and
PICARD~\citep{scholak2021picard} achieve competitive accuracy by training
a task-specific model on Spider data; RESDSQL attains EX\,=\,0.791 with a
T5-3B encoder-decoder, while PICARD uses constrained decoding to reach 0.754.
Both require retraining when the target domain changes, limiting their
use in cross-domain edge deployment.

Prompting-based approaches operate without fine-tuning by leveraging
few-shot examples and chain-of-thought decomposition.
DAIL-SQL~\citep{gao2023dailsql} achieves EX\,=\,0.866 by constructing
prompts with schema DDL, sample rows, and structurally similar examples
selected via sql2skeleton similarity.
DIN-SQL~\citep{pourreza2023dinsql} decomposes queries into subproblems,
achieving EX\,=\,0.826.
DCG-SQL~\citep{lee2025dcgsql} further enhances in-context learning via
a deep contextual schema link graph, attaining 0.876 on Spider~1.0.
These approaches require proprietary cloud models and cannot be
reproduced on consumer hardware.

Code-specialised local models such as Qwen2.5-Coder~\citep{hui2024qwen25coder}
enable Text-to-SQL on edge hardware, but systematic ablation studies of
pipeline components under resource constraints remain scarce in the
literature~\citep{liu2025survey}.

\subsection{Quantization for Neural Retrieval}

\citet{jacob2018quantization} established the theoretical framework for
INT8 inference: weights are stored as 8-bit integers with a per-tensor scale
factor, and activations are quantized dynamically on a per-token basis during
the forward pass.
This approach requires no calibration data and applies post-training.
Q8BERT~\citep{zafrir2019q8bert} demonstrated that 8-bit quantization of all
linear layers in BERT retains task accuracy on GLUE, establishing INT8 as a
reliable default for encoder-only transformers.
Q-BERT~\citep{shen2020qbert} subsequently showed that Hessian-based analysis
identifies the most sensitive layers, enabling mixed-precision quantization
with smaller accuracy loss.

For large language models, LLM.int8()~\citep{dettmers2022llmint8} isolates
outlier activations in FP16 while quantizing the remainder to INT8,
enabling 8-bit inference at billion-parameter scale.
SmoothQuant~\citep{xiao2023smoothquant} achieves accurate INT8
weight--activation quantization by migrating the quantization difficulty from
activations to weights via a per-channel scaling transformation.
GPTQ~\citep{dettmers2022gptq} and AWQ~\citep{lin2023awq} extend
post-training quantization to 4-bit precision with minimal accuracy loss.

Despite this extensive literature, the impact of INT8 quantization on
dense retrieval encoders in Text-to-SQL pipelines has not been studied.
The present work fills this gap by conducting the first
per-layer and per-sublayer sensitivity analysis for a schema-linking
retriever, measuring both retrieval quality and downstream execution accuracy.

\subsection{Schema Retrieval and Linking}

Schema linking---identifying which tables and columns in a database are
relevant to a natural-language question---is a prerequisite for accurate
cross-database SQL generation~\citep{lei2020schema}.
\citet{wang2020ratsql} introduced relation-aware schema encoding via a
typed graph representation, establishing the importance of structural
relationships between question tokens and schema elements.
Dense retrieval with sentence embeddings~\citep{reimers2019sbert}
subsequently outperformed sparse retrieval for schema linking.

Recent work has scaled schema linking to thousands of tables.
RASL~\citep{eben2025rasl} proposes retrieval-augmented schema linking for
massive databases, while LinkAlign~\citep{wang2025linkalign} addresses
multi-database settings where the correct database must be identified before
table retrieval.
\citet{nahid2025rethinking} introduces a context-aware bidirectional retrieval
approach that re-ranks candidates using both question-to-schema and
schema-to-question signals.
AutoLink~\citep{wang2025autolink} extends schema exploration to unseen schemas
without manual annotation.

The Snowflake arctic-embed-m model~\citep{merrick2024arctic} is an asymmetric
embedding model trained for passage retrieval that generalises well
across domains without fine-tuning, making it suitable for the
cross-database schema retrieval task.
